% SoftwareX Original Software Publication draft.
\documentclass[preprint,12pt]{elsarticle}
\usepackage[T1]{fontenc}
\usepackage{lmodern}
\usepackage{microtype}
\usepackage{amsmath}
\usepackage{booktabs}
\usepackage{hyperref}
\usepackage{lineno}
\usepackage{tabularx}
\hypersetup{colorlinks=true,allcolors=blue}
\newcolumntype{Y}{>{\raggedright\arraybackslash}X}
\setlength{\parindent}{0pt}

\journal{SoftwareX}

\begin{document}
\begin{frontmatter}

\title{RosettaX: reproducible calibration workflows for flow-cytometry FCS data}

\author[amsterdamumc]{Martin Poinsinet de Sivry-Houle\corref{corresponding}}
\ead{martin.poinsinet.de.sivry@gmail.com}
\cortext[corresponding]{Corresponding author}
\affiliation[amsterdamumc]{%
  organization={Amsterdam UMC},
  city={Amsterdam},
  country={the Netherlands}}

\begin{abstract}
RosettaX is an open-source local application for calibrating exported
flow-cytometry data in the Flow Cytometry Standard (FCS) format. It guides
users through histogram and peak review, reference-table approval, fitting,
and application of saved calibrations. Fluorescence workflows fit an
empirical power law between detector values and assigned reference values.
Scattering workflows fit measured peaks against optical coupling calculated
with a configurable Mie-scattering and detector-geometry model. Each
calibration is saved as structured JSON containing the fit, source channel,
reference inputs, and workflow-specific context. Applying a calibration
exports calibrated FCS data with a PDF report, keeping the result and its
provenance together for review and reuse.
\end{abstract}

\begin{keyword}
flow cytometry \sep calibration \sep fluorescence \sep light scattering \sep
FCS \sep Mie theory
\end{keyword}
\end{frontmatter}

\modulolinenumbers[1]
\linenumbers

\begin{table}[t]
\centering
\footnotesize
\caption{Code metadata.}
\label{tab:code-metadata}
\begin{tabularx}{\textwidth}{@{}p{0.35in}p{2.25in}Y@{}}
\toprule
Nr. & Code metadata description & Metadata \\
\midrule
C1 & Current code version & 0.9.0 \\
C2 & Permanent link to code/repository used for this version &
\url{https://github.com/MartinPdeS/RosettaX/tree/v0.9.0} \\
C3 & Permanent link to reproducible capsule & Not available; tagged source,
documentation, and tests are versioned with the repository. \\
C4 & Legal code license & GNU Affero General Public License v3.0 or later \\
C5 & Code versioning system used & Git \\
C6 & Software code languages, tools, and services used & Python, Dash, NumPy,
Plotly, PyMieSim, and PostgreSQL-compatible persistence \\
C7 & Compilation requirements, operating environments, and dependencies &
Python 3.10 or later; dependencies declared in \texttt{pyproject.toml} \\
C8 & Link to developer documentation/manual &
\url{https://martinpdes.github.io/RosettaX/} \\
C9 & Support email for questions &
\href{mailto:martin.poinsinet.de.sivry@gmail.com}{martin.poinsinet.de.sivry@gmail.com} \\
\bottomrule
\end{tabularx}
\end{table}

\begin{table}[t]
\centering
\footnotesize
\caption{Executable software metadata.}
\label{tab:software-metadata}
\begin{tabularx}{\textwidth}{@{}p{0.35in}p{2.25in}Y@{}}
\toprule
Nr. & Software metadata description & Metadata \\
\midrule
S1 & Current software version & 0.9.0 \\
S2 & Permanent link to executables of this version &
\url{https://github.com/MartinPdeS/RosettaX/releases/tag/v0.9.0} \\
S3 & Permanent link to reproducible capsule & Not available; use the tagged
source and release artifacts. \\
S4 & Legal software license & GNU Affero General Public License v3.0 or later \\
S5 & Computing platforms/operating systems & Linux, macOS, and Microsoft Windows \\
S6 & Installation requirements and dependencies & Python 3.10 or later for
source installation; release bundles are supplied for supported platforms. \\
S7 & Link to user manual & \url{https://martinpdes.github.io/RosettaX/} \\
S8 & Support email for questions &
\href{mailto:martin.poinsinet.de.sivry@gmail.com}{martin.poinsinet.de.sivry@gmail.com} \\
\bottomrule
\end{tabularx}
\end{table}

\section{Motivation and significance}

Flow-cytometry instruments export event-level measurements whose detector
scales depend on the instrument and acquisition context. A reusable
calibration must therefore retain more than fitted coefficients: the selected
channel, approved reference values, optical assumptions where applicable, and
the context in which the transformation was produced. FCS provides a standard
format for exchange of these event data \cite{Spidlen2010}, but it does not
itself encode a calibration workflow.

RosettaX provides a local, guided workflow for this post-acquisition step. It
operates on exported FCS data rather than controlling an instrument. The
application helps users inspect distributions, identify and review peaks, pair
them with reference values or modeled standards, save a structured calibration,
and apply that record to experimental files. Its contribution is to place
fluorescence and scattering calibration, saved context, and export provenance
in one workflow without representing processed FCS values as raw optical
measurements.

\section{Software description}

\subsection{Calibration workflows}

For fluorescence, RosettaX fits approved positive detector-reference pairs in
logarithmic space:
\begin{equation}
  \log_{10}(y) = m\log_{10}(x)+b,\qquad y=10^b x^m.
\end{equation}
The saved record contains the slope, intercept, prefactor, $R^2$, point count,
and the reference pairs that supported the fit. This makes the fitted
transformation reviewable without reopening the original calibration session.

For scattering, RosettaX calculates the coupling of calibration standards with
PyMieSim using a Gaussian source and photodiode detector model. The user can
define wavelength, material refractive indices, and detector geometry,
including numerical apertures and angular settings. The instrument response
relates the measured peak to modeled coupling, by default with a
zero-intercept linear fit. The physical model follows Mie scattering for
spherical particles \cite{Mie1908} and uses the versioned PyMieSim dependency
\cite{PyMieSim}; it does not assume that scatter peak is directly proportional
to particle diameter.

\subsection{Saved records and application safeguards}

Both workflows use a common JSON wrapper identifying the calibration kind and
creation context. Fluorescence payloads preserve their empirical fit and
reference points. Scattering payloads additionally preserve the
instrument-response fit, modeled calibration-standard relation, and optical
metadata. RosettaX validates the record and source channel before export.
When a scattering target relation is not strictly monotonic, it selects the
largest monotonic branch and reports that decision rather than silently treating
the inverse as single valued.

\section{Illustrative use}

A typical fluorescence calibration begins with a bead FCS file. The user
selects a detector channel, reviews the visible populations, approves the
peak-reference table, and examines the power-law fit before saving the JSON
record. The same record can later be selected with one or more experimental
FCS files to add a calibrated output column.

For scattering, the analogous workflow combines measured calibration-standard
peaks with a detector preset or custom collection geometry and material
settings. RosettaX models the standard relation, fits the instrument response,
and stores the resulting optical context. On application, the software
converts measured values to estimated coupling and then evaluates the selected
target relation. The resulting FCS artifacts may be bundled with a PDF report
that records selected calibration, input files, output columns, and warnings.

\section{Impact}

RosettaX is intended for laboratories that need to make their calibration
decisions inspectable and reusable across exported FCS files. Its JSON record
separates calibration families while retaining the evidence required to
interpret each one, and its co-exported report captures what occurred during a
specific application run. The application supports detector presets for
selected BD Biosciences, Beckman Coulter, Apogee, and generic instrument
configurations while preserving custom optical settings in the calibration
record. It is not a replacement for instrument control, a universal optical
model, or scientific judgment about reference materials and target assumptions.

\section{Conclusions}

RosettaX supplies a documented local route from exported FCS measurements to
reviewable fluorescence or scattering calibrations and provenance-bearing
exports. By storing both fitted relations and their interpretation context, it
helps users reuse calibrations while retaining the information needed to assess
them. The tagged source, documentation, and software archive are available
openly \cite{RosettaX}.

\section*{Data availability}

No new experimental data were generated for this software paper. RosettaX
operates on user-supplied FCS files; users should share data only in accordance
with their ethical, privacy, and institutional requirements.

\section*{Declaration of competing interest}

The author declaration requires confirmation before submission.

\section*{Acknowledgements}

Funding and acknowledgement statements require confirmation before submission.

\section*{Generative-AI disclosure}

The authors must replace this statement before submission with an accurate
description of any generative-AI use and their review of the final manuscript.

\bibliographystyle{elsarticle-num}
\bibliography{references}

\end{document}
